% Options for packages loaded elsewhere
\PassOptionsToPackage{unicode}{hyperref}
\PassOptionsToPackage{hyphens}{url}
\PassOptionsToPackage{dvipsnames,svgnames,x11names}{xcolor}
%
\documentclass[
  10pt,
  a4paper,
  DIV=11]{scrartcl}

\usepackage{amsmath,amssymb}
\usepackage{iftex}
\ifPDFTeX
  \usepackage[T1]{fontenc}
  \usepackage[utf8]{inputenc}
  \usepackage{textcomp} % provide euro and other symbols
\else % if luatex or xetex
  \usepackage{unicode-math}
  \defaultfontfeatures{Scale=MatchLowercase}
  \defaultfontfeatures[\rmfamily]{Ligatures=TeX,Scale=1}
\fi
\usepackage{lmodern}
\ifPDFTeX\else  
    % xetex/luatex font selection
  \setmainfont[]{Arial}
\fi
% Use upquote if available, for straight quotes in verbatim environments
\IfFileExists{upquote.sty}{\usepackage{upquote}}{}
\IfFileExists{microtype.sty}{% use microtype if available
  \usepackage[]{microtype}
  \UseMicrotypeSet[protrusion]{basicmath} % disable protrusion for tt fonts
}{}
\makeatletter
\@ifundefined{KOMAClassName}{% if non-KOMA class
  \IfFileExists{parskip.sty}{%
    \usepackage{parskip}
  }{% else
    \setlength{\parindent}{0pt}
    \setlength{\parskip}{6pt plus 2pt minus 1pt}}
}{% if KOMA class
  \KOMAoptions{parskip=half}}
\makeatother
\usepackage{xcolor}
\usepackage[top=30mm,left=20mm,heightrounded]{geometry}
\setlength{\emergencystretch}{3em} % prevent overfull lines
\setcounter{secnumdepth}{3}
% Make \paragraph and \subparagraph free-standing
\ifx\paragraph\undefined\else
  \let\oldparagraph\paragraph
  \renewcommand{\paragraph}[1]{\oldparagraph{#1}\mbox{}}
\fi
\ifx\subparagraph\undefined\else
  \let\oldsubparagraph\subparagraph
  \renewcommand{\subparagraph}[1]{\oldsubparagraph{#1}\mbox{}}
\fi


\providecommand{\tightlist}{%
  \setlength{\itemsep}{0pt}\setlength{\parskip}{0pt}}\usepackage{longtable,booktabs,array}
\usepackage{calc} % for calculating minipage widths
% Correct order of tables after \paragraph or \subparagraph
\usepackage{etoolbox}
\makeatletter
\patchcmd\longtable{\par}{\if@noskipsec\mbox{}\fi\par}{}{}
\makeatother
% Allow footnotes in longtable head/foot
\IfFileExists{footnotehyper.sty}{\usepackage{footnotehyper}}{\usepackage{footnote}}
\makesavenoteenv{longtable}
\usepackage{graphicx}
\makeatletter
\def\maxwidth{\ifdim\Gin@nat@width>\linewidth\linewidth\else\Gin@nat@width\fi}
\def\maxheight{\ifdim\Gin@nat@height>\textheight\textheight\else\Gin@nat@height\fi}
\makeatother
% Scale images if necessary, so that they will not overflow the page
% margins by default, and it is still possible to overwrite the defaults
% using explicit options in \includegraphics[width, height, ...]{}
\setkeys{Gin}{width=\maxwidth,height=\maxheight,keepaspectratio}
% Set default figure placement to htbp
\makeatletter
\def\fps@figure{htbp}
\makeatother

%% load packages
\usepackage{sectsty}
\usepackage{fontspec}
\usepackage{scrlayer-scrpage}
\usepackage{fontspec}

% Define title formatting

%% Set font sizes for sections and subsections
\sectionfont{\fontsize{11}{16}\selectfont\bfseries} % Arial 14pt bold
\subsectionfont{\fontsize{10}{14}\selectfont} % Arial 12pt

\lehead[test1]{\pagemark}
\cehead[]{}
\rehead[test 2]{\leftmark}
\lohead[{\fontsize{8}{10}\selectfont\upshape\textbf{Institut Sekundarstufe I}} \\
\fontsize{8}{10}\selectfont\upshape Fabrikstrasse 8, CH-3012 Bern \\
\fontsize{8}{10}\selectfont\upshape T +41 31 309 21 15, contactdesk@phbern.ch, www.phbern.ch
]{\rightmark}
\rohead[PHBern]{}
\KOMAoption{captions}{tableheading}
\makeatletter
\@ifpackageloaded{tcolorbox}{}{\usepackage[skins,breakable]{tcolorbox}}
\@ifpackageloaded{fontawesome5}{}{\usepackage{fontawesome5}}
\definecolor{quarto-callout-color}{HTML}{909090}
\definecolor{quarto-callout-note-color}{HTML}{0758E5}
\definecolor{quarto-callout-important-color}{HTML}{CC1914}
\definecolor{quarto-callout-warning-color}{HTML}{EB9113}
\definecolor{quarto-callout-tip-color}{HTML}{00A047}
\definecolor{quarto-callout-caution-color}{HTML}{FC5300}
\definecolor{quarto-callout-color-frame}{HTML}{acacac}
\definecolor{quarto-callout-note-color-frame}{HTML}{4582ec}
\definecolor{quarto-callout-important-color-frame}{HTML}{d9534f}
\definecolor{quarto-callout-warning-color-frame}{HTML}{f0ad4e}
\definecolor{quarto-callout-tip-color-frame}{HTML}{02b875}
\definecolor{quarto-callout-caution-color-frame}{HTML}{fd7e14}
\makeatother
\makeatletter
\@ifpackageloaded{caption}{}{\usepackage{caption}}
\AtBeginDocument{%
\ifdefined\contentsname
  \renewcommand*\contentsname{Inhaltsverzeichnis}
\else
  \newcommand\contentsname{Inhaltsverzeichnis}
\fi
\ifdefined\listfigurename
  \renewcommand*\listfigurename{Abbildungsverzeichnis}
\else
  \newcommand\listfigurename{Abbildungsverzeichnis}
\fi
\ifdefined\listtablename
  \renewcommand*\listtablename{Tabellenverzeichnis}
\else
  \newcommand\listtablename{Tabellenverzeichnis}
\fi
\ifdefined\figurename
  \renewcommand*\figurename{Abbildung}
\else
  \newcommand\figurename{Abbildung}
\fi
\ifdefined\tablename
  \renewcommand*\tablename{Tabelle}
\else
  \newcommand\tablename{Tabelle}
\fi
}
\@ifpackageloaded{float}{}{\usepackage{float}}
\floatstyle{ruled}
\@ifundefined{c@chapter}{\newfloat{codelisting}{h}{lop}}{\newfloat{codelisting}{h}{lop}[chapter]}
\floatname{codelisting}{Listing}
\newcommand*\listoflistings{\listof{codelisting}{Listingverzeichnis}}
\usepackage{amsthm}
\theoremstyle{definition}
\newtheorem{exercise}{Übungsaufgabe}[section]
\theoremstyle{remark}
\AtBeginDocument{\renewcommand*{\proofname}{Beweis}}
\newtheorem*{remark}{Anmerkung}
\newtheorem*{solution}{Lösung}
\newtheorem{refremark}{Anmerkung}[section]
\newtheorem{refsolution}{Lösung}[section]
\makeatother
\makeatletter
\makeatother
\makeatletter
\@ifpackageloaded{caption}{}{\usepackage{caption}}
\@ifpackageloaded{subcaption}{}{\usepackage{subcaption}}
\makeatother
\ifLuaTeX
\usepackage[bidi=basic]{babel}
\else
\usepackage[bidi=default]{babel}
\fi
\babelprovide[main,import]{ngerman}
\ifPDFTeX
\else
\babelfont{rm}[]{Arial}
\fi
% get rid of language-specific shorthands (see #6817):
\let\LanguageShortHands\languageshorthands
\def\languageshorthands#1{}
\ifLuaTeX
  \usepackage{selnolig}  % disable illegal ligatures
\fi
\usepackage{bookmark}

\IfFileExists{xurl.sty}{\usepackage{xurl}}{} % add URL line breaks if available
\urlstyle{same} % disable monospaced font for URLs
\hypersetup{
  pdftitle={Funktionen},
  pdfauthor={Richard Conrardy},
  pdflang={de},
  colorlinks=true,
  linkcolor={(3,1.5)},
  filecolor={Maroon},
  citecolor={Blue},
  urlcolor={Blue},
  pdfcreator={LaTeX via pandoc}}

\title{Funktionen}
\usepackage{etoolbox}
\makeatletter
\providecommand{\subtitle}[1]{% add subtitle to \maketitle
  \apptocmd{\@title}{\par {\large #1 \par}}{}{}
}
\makeatother
\subtitle{Einführung ins Konzept}
\author{Lynn Gurtner}
\date{21.09.2024}

\begin{document}
\maketitle

\renewcommand*\contentsname{Inhaltsverzeichnis}
{
\hypersetup{linkcolor=}
\setcounter{tocdepth}{3}
\tableofcontents
}
\section{Lektion 1}\label{lektion-1}

\subsection{Organisation}\label{organisation}

\begin{longtable}[]{@{}
  >{\raggedright\arraybackslash}p{(\columnwidth - 8\tabcolsep) * \real{0.0385}}
  >{\raggedright\arraybackslash}p{(\columnwidth - 8\tabcolsep) * \real{0.2500}}
  >{\raggedright\arraybackslash}p{(\columnwidth - 8\tabcolsep) * \real{0.3205}}
  >{\raggedright\arraybackslash}p{(\columnwidth - 8\tabcolsep) * \real{0.0609}}
  >{\raggedright\arraybackslash}p{(\columnwidth - 8\tabcolsep) * \real{0.3301}}@{}}
\caption{Lektionsplan Lektion 1}\tabularnewline
\toprule\noalign{}
\begin{minipage}[b]{\linewidth}\raggedright
Zeitangabe
\end{minipage} & \begin{minipage}[b]{\linewidth}\raggedright
Aufgabe SuS
\end{minipage} & \begin{minipage}[b]{\linewidth}\raggedright
Aufgabe LP
\end{minipage} & \begin{minipage}[b]{\linewidth}\raggedright
LUKAS
\end{minipage} & \begin{minipage}[b]{\linewidth}\raggedright
Methode
\end{minipage} \\
\midrule\noalign{}
\endfirsthead
\toprule\noalign{}
\begin{minipage}[b]{\linewidth}\raggedright
Zeitangabe
\end{minipage} & \begin{minipage}[b]{\linewidth}\raggedright
Aufgabe SuS
\end{minipage} & \begin{minipage}[b]{\linewidth}\raggedright
Aufgabe LP
\end{minipage} & \begin{minipage}[b]{\linewidth}\raggedright
LUKAS
\end{minipage} & \begin{minipage}[b]{\linewidth}\raggedright
Methode
\end{minipage} \\
\midrule\noalign{}
\endhead
\bottomrule\noalign{}
\endlastfoot
5 min & Materialien bereitlegen & Begrüssung, Ankündigung neues Thema &
Ankommen & Frontalunterricht \\
15 min & Aufmerksam zuhören, Fragen stellen bei Unklarheiten. Beispiele
notieren und gemeinsam lösen. & Einführung Funktionen Teil 1 mit Skript
als Orientierung. Mehrere Beispiele werden an der Wandtafel gezeigt,
Begriffe werden erklärt. Zugriff zum Skript erst am Ende der Erklärung.
& Konfrontation & An der Wandtafel und Beamer, Frontalunterricht \\
10 min & Durchlesen des Skripts. Teil 1. Unklarheiten mit LP klären,
Wichtige Begriffe markieren & Zugriff zum Skript erteilen. Fragen der
SuS klären. & Erarbeitungsaufgabe & Skript online hochladen, SuS machen
Einzelarbeit beim Durchlesen, Fragen mit LP oder mit Nachbar*in
klären. \\
15 min & Aufgaben 1-4 lösen & Bereit für Fragen, allenfalls erneute
Erklärungen an der Wandtafel. & Übungsaufgabe & Aufgaben können in
Einzel- oder Partnerarbeit gelöst werden. Entweder direktes Lösen auf
dem Tablet oder Übersetzen auf Papier. \\
\end{longtable}

\subsection{Skript}\label{skript}

Skript Einführung Funktionen Teil 1 Was ist eine Funktion?

Die Funktion ist eines der wichtigsten «Objekte» der Mathematik. In der
Biologie oder im Physik werden sie oftmals auch als \textbf{Regeln,
Gesetze oder Formeln} bezeichnet. Wir stellen uns Funktionen in der
Mathematik als \textbf{kleine Maschinen} vor. In diese Maschine wird ein
\textbf{Input} gegeben, dieser wird dann nach einer genauen Regel
verarbeitet und als neuer \textbf{Output} wieder aus der Maschine
gegeben.
\href{https://v-fe.lehrplan.ch/101DeKLWBVL4Rk2n6hmVf3GLKffyRqX4Y}{MA.3.A.1.m}

Eine Funktion wird üblicherweise mit dem Symbol f bezeichnet. Wobei der
Input mit einem x und der Output mit einem y

gekennzeichnet wird. Sind aber mehrere verschiedene Funktionen
vorhanden, werden diese mit weiteren Buchstaben bezeichnet.

Hängen wir diese Informationen nun zusammen, erhalten wir eine Maschine:

\$ \text{Input } x \$

\[
\downarrow
\]

\[
\boxed{f} \quad \text{Regel: } f(x) = x + 3
\]

\[
\downarrow
\]

\$ \text{Output } y \$

Der Input wird also in die Maschine gegeben, wird dann in der Maschine
verändert und kommt als neuer Output wieder heraus. Die Regel dieser
Funktion lautet x+3 . Das heisst, addiere den Input mit 3. Wenn also der
Input 2 ist, so ist der Output 2+3=5.

Hier ist eine Maschine \(g\),

die mit einer anderen Regel funktioniert.

\$ \text{Input } x \$

\[
\downarrow
\]

\[
\boxed{g} \quad \text{Regel: } g(x) = x^2
\]

\[
\downarrow
\]

\$ \text{Output } y \$

Bei dieser Funktion wird der Inhalt quadriert. Wenn also \(x = 4\), so
ist \(y\), oder der Output, \(4 \cdot 4 = 16\).

Nun gibt es verschiedene \textbf{Notationen}, wie man eine Funktion ohne
die Maschine beschreiben kann. Wir werden uns mit der meistgebrauchten
Notation beschäftigen.

Diese sieht folgendermassen aus: \(ƒ(x)\) \textbf{= Regel der Funktion}.

Damit die Funktion berechnet werden kann, wird das \(x\) definiert,
sodass man den \(y\)-Wert finden kann. Andererseits kann man auch den
\textbf{\(y\)-Wert} definieren und so mit Hilfe der Funktionsregel auf
den \textbf{\(x\)-Wert} schliessen.

Wenn wir nun unsere Maschinen von oben betrachten, sehen diese mit der
Notation so aus:

\[ f(x) = x + 3 \]

\[ g(x) = x^2 \]

\textbf{Löse Aufgaben 1 \& 2 (Grundanforderung)}

\textbf{Aufgabe 3:} obligatorisch für Sek und Spez-Sek, freiwillig für
Real

\textbf{Aufgabe 4:} obligatorisch für Spez-Sek, freiwillig für Sek und
Real

\subsection{Teil 2 Eigenschaften von Funktionen: Welche Gleichung ist
keine
Funktion?}\label{teil-2-eigenschaften-von-funktionen-welche-gleichung-ist-keine-funktion}

Für eine Funktion kann nicht einfach eine willkürliche Gleichung gewählt
werden. Es gibt gewisse Regeln, die beachtet werden müssen. Um diese
etwas anschaulicher zu gestalten, benutzen wir dazu die \textbf{Mengen}.

Jede Funktion besteht aus 2 Mengen. Diese Mengen werden als
\textbf{Definitionsbereich} und als \textbf{Wertebereich} bezeichnet.

Der Definitionsbereich einer Funktion gibt an welche Werte man in die
Funktion einsetzen darf. Er beschreibt also alle Inputs oder x-Werte.

Der Wertebereich beschreibt alle möglichen Ergebnisse, die aus der
Maschine ausgespuckt werden können. Es sind also die Outputs oder
y-Werte einer Funktion. Die Zahlen des Wertebereichs entstehen also
durch die Zahlen des Definitionsbereichs.

Damit eine Funktion nun gültig ist, kann \textbf{jeder Input nur einen
einzigen Output} besitzen. Das bedeutet jedem Element im
Definitionsbereich wird genau 1 Element im Wertebereich zugeteilt.
Sobald die Gleichung mehrere Lösungen hat, ist es keine
Funktionsgleichung.

Lasst uns das in einem Bild veranschaulichen:

\includegraphics{.pdf}

Wie man sehen kann, verändern sich die Formen oder auch Zahlen durch die
Funktionsregeln, jedoch wird jedem Input genau 1 Output zugestellt.
Somit sind die Funktionen gültig.

\subsubsection{Ist die folgende Gleichung eine Funktionsgleichung?
Begründe:}\label{ist-die-folgende-gleichung-eine-funktionsgleichung-begruxfcnde}

\[ f(2) = x^2 - 4 \]

Ebenso wird die Gleichung nicht als gültig anerkannt, wenn sie kein
Ergebnis ausspuckt.

Ein Beispiel für eine Gleichung mit ungültigem Output ist:

\[ f(3) = \frac{1}{x - 3} \]

Diese Gleichung ist ungültig, da die Division durch Null nicht erlaubt
ist. Die Gleichung hat also keinen gültigen Output.

\subsubsection{MA.3.B.1.i}\label{ma.3.b.1.i}

Löse Aufgaben 5-9. Aufgabe 9 = freiwillige Erweiterung.

\subsubsection{Wichtige Erkenntnisse und
Begriffe:}\label{wichtige-erkenntnisse-und-begriffe}

\begin{itemize}
\tightlist
\item
  Eine Funktion kann als kleine Maschine betrachtet werden.
\item
  Eine Funktion wird mit dem Symbol ƒ bezeichnet.
\item
  \textbf{Input} = x-Wert der Funktion.
\item
  \textbf{Output} = y-Wert der Funktion.
\item
  Jede Funktion besitzt eine Regel, nach der der Output berechnet wird.
\item
  \[ ƒ(x) \] = Regel der Funktion ist die Notation einer Funktion.
\item
  Der \textbf{Definitionsbereich} beschreibt die Inputs.
\item
  Der \textbf{Wertebereich} beschreibt die Outputs.
\end{itemize}

Wenn du noch eine zusätzliche Erklärung zum Thema Funktionen brauchst,
kannst du gerne noch das Video unter folgendem Link anschauen:

\href{https://youtu.be/myLx0d5wmHw?si=vB-aI0zptBWo8Rt4}{Video zur
Erklärung der Funktionen}

\textbf{Achtung:} Hier werden bereits die Funktionsgraphen erklärt,
welche wir uns in einem nächsten Schritt anschauen. Du darfst dich aber
gerne schon darüber informieren.

\section{Lektion 2}\label{lektion-2}

\begin{longtable}[]{@{}
  >{\raggedright\arraybackslash}p{(\columnwidth - 8\tabcolsep) * \real{0.0241}}
  >{\raggedright\arraybackslash}p{(\columnwidth - 8\tabcolsep) * \real{0.3320}}
  >{\raggedright\arraybackslash}p{(\columnwidth - 8\tabcolsep) * \real{0.3843}}
  >{\raggedright\arraybackslash}p{(\columnwidth - 8\tabcolsep) * \real{0.0523}}
  >{\raggedright\arraybackslash}p{(\columnwidth - 8\tabcolsep) * \real{0.2072}}@{}}
\caption{Lektionsplan Lektion 2}\tabularnewline
\toprule\noalign{}
\begin{minipage}[b]{\linewidth}\raggedright
Zeitangabe
\end{minipage} & \begin{minipage}[b]{\linewidth}\raggedright
Aufgabe SuS
\end{minipage} & \begin{minipage}[b]{\linewidth}\raggedright
Aufgabe LP
\end{minipage} & \begin{minipage}[b]{\linewidth}\raggedright
LUKAS
\end{minipage} & \begin{minipage}[b]{\linewidth}\raggedright
Methode
\end{minipage} \\
\midrule\noalign{}
\endfirsthead
\toprule\noalign{}
\begin{minipage}[b]{\linewidth}\raggedright
Zeitangabe
\end{minipage} & \begin{minipage}[b]{\linewidth}\raggedright
Aufgabe SuS
\end{minipage} & \begin{minipage}[b]{\linewidth}\raggedright
Aufgabe LP
\end{minipage} & \begin{minipage}[b]{\linewidth}\raggedright
LUKAS
\end{minipage} & \begin{minipage}[b]{\linewidth}\raggedright
Methode
\end{minipage} \\
\midrule\noalign{}
\endhead
\bottomrule\noalign{}
\endlastfoot
15 min & Zuhören, Fragen stellen bei Unklarheiten & Einführung des 2.
Teils (Mengen). Unklarheiten wiederholen, Begriffe erklären mit
Beispielen (Erklärung mit Menge = Beispiel 1 bei Beispielen einer
Funktion von Serlo als Hilfe) & Konfrontation und Reflexion &
Frontalunterricht, Wandtafel und Beamer \\
20 min & Aufgaben 5-9 lösen. Bei Unklarheiten LP fragen oder verlinkte
Videos anschauen. & SuS unterstützen beim Lösen der Aufgaben. Allenfalls
Aufgaben im Plenum erklären. Einige schwierigere Aufgaben im
Frontalunterricht lösen, falls erwünscht von den SuS. & Übungs- und
Vertiefungsaufgaben & Einzel- oder Partnerarbeit. SuS arbeiten am Platz
mit Tablets. \\
10 min & Letzter Teil: Wichtige Begriffe und Erkenntnisse durchlesen.
Unklarheiten melden. & Wiederholung der wichtigen Begriffe und
Erkenntnisse. Schwierige Aufgaben besprechen und gemeinsam lösen.
Allenfalls von SuS erklären lassen. & Reflexion und Synthese & Zuerst
Einzelarbeit beim Lesen, danach Besprechung und Klärung im
Frontalunterricht \\
\end{longtable}

\section{Übungen
Funktionsgleichungen}\label{uxfcbungen-funktionsgleichungen}

\begin{exercise}[]\protect\hypertarget{exr-line}{}\label{exr-line}

Finde den y-Wert

\begin{flalign*}
&\text{a.} \quad &f(2)  &=  \ 3x + 2 \\
&\text{b.} \quad &g(7)  &=  \ 2x - 1 \\
&\text{c.} \quad &h(8)  &=  \ 12x + 3 \\
&\text{d.} \quad &i(1)  &=  \ 12x \cdot 3 - 9 \\
&\text{e.} \quad &k(2)  &=  \ 2x
\end{flalign*}

\end{exercise}

\begin{exercise}[]\protect\hypertarget{exr-line}{}\label{exr-line}

Erstelle eine Funktionsgleichung mit

\begin{enumerate}
\def\labelenumi{\alph{enumi}.}
\tightlist
\item
  \(x = 2\)
\item
  \(x = −2\)
\item
  \(x = 16\)
\item
  \(x = 1.5\)
\item
  \(x = 12\)
\item
  \(x = 7 − 4 / 2\)
\end{enumerate}

\end{exercise}

\begin{exercise}[]\protect\hypertarget{exr-line}{}\label{exr-line}

Betrachte die folgenden, mit Worten beschriebenen Funktionen. Zeichne
für jede Funktion die Maschinen-Darstellung und schreibe auch die
Funktionsgleichung auf. (MA.1.A.4.j)

\begin{enumerate}
\def\labelenumi{\arabic{enumi}.}
\tightlist
\item
  Der Input wird durch 2 geteilt, und das Ergebnis wird um 1 erhöht.
\item
  Der Input wird um 2.5 vermindert, und das Ergebnis wird mit 10
  multipliziert.
\item
  Der Input wird hoch drei gerechnet, und das Ergebnis wird mit 2
  multipliziert.
\item
  Der Input wird mit 3 multipliziert, das Ergebnis wird hoch drei
  gerechnet.
\end{enumerate}

\end{exercise}

\begin{exercise}[]\protect\hypertarget{exr-line}{}\label{exr-line}

Bestimme den Input für den gegebenen Output

\subsubsection{1.}\label{section}

x = ?

↓ ⌈ƒ⌋\\
Regel: f(x) = 3x − 1\\
↓\\
y = 0

\subsubsection{2.}\label{section-1}

x = ?

↓ ⌈ƒ⌋\\
Regel: f(x) = 2x\\
↓\\
y = 36

\subsubsection{3.}\label{section-2}

x = ?

↓ ⌈ƒ⌋\\
Regel: f(x) = (x − 1)²\\
↓\\
y = 16

\subsubsection{4.}\label{section-3}

x = ?

↓ ⌈ƒ⌋ Regel: f(x) = 1x\\
↓\\
y = 0

\end{exercise}

\begin{exercise}[]\protect\hypertarget{exr-line}{}\label{exr-line}

f ist die Funktion, die jeder Frucht in der Menge die richtige Farbe
zuordnet.

Zeichne die Mengen des Definitionsbereichs und des Wertebereichs.
Verbinde die Elemente durch Pfeile (wie im Skript) und notiere die
Funktionsgleichungen für:

\begin{itemize}
\tightlist
\item
  Apfel\\
\item
  Banane\\
\item
  Birne\\
\item
  Mango\\
\item
  Pflaume\\
\item
  Tomate
\end{itemize}

Funktionsgleichungen können auch im Alltag nützlich sein. Schauen wir
uns eine solche Situation etwas genauer an. (MA.3.C.3.g)

\end{exercise}

\begin{exercise}[]\protect\hypertarget{exr-line}{}\label{exr-line}

Die Fahrzeit berechnen mit Hilfe einer Funktionsgleichung

Für die Berechnung der Fahrzeit benötigen wir 2 verschiedene Angaben:

\begin{enumerate}
\def\labelenumi{\arabic{enumi}.}
\tightlist
\item
  Die Strecke\\
\item
  Die Geschwindigkeit
\end{enumerate}

Für die Strecke benutzen wir den Buchstaben \textbf{d} (distance) mit
der Maßeinheit km und für die Geschwindigkeit den Buchstaben \textbf{v}
(velocity) mit der Maßeinheit km/h.\\
Wenn wir nun die Fahrzeit berechnen wollen, benötigen wir die folgende
Gleichung:

\textbf{f(t) = d / v}

Die Strecke geteilt durch die Geschwindigkeit ergibt also die Fahrzeit.

\end{exercise}

\begin{exercise}[]\protect\hypertarget{exr-line}{}\label{exr-line}

Berechne die Fahrzeit. Gib in Minuten an

\begin{enumerate}
\def\labelenumi{\arabic{enumi}.}
\tightlist
\item
  t = 120, v = 60
\item
  t = 30, v = 60
\item
  t = 100, v = 80
\item
  t = 70, v = 40
\end{enumerate}

\end{exercise}

\begin{exercise}[]\protect\hypertarget{exr-line}{}\label{exr-line}

Berechne die Fahrzeit von der Schule zu deinem Zuhause und an deinen
Lieblings-Ferienort mit

\begin{enumerate}
\def\labelenumi{\arabic{enumi}.}
\tightlist
\item
  v = 80
\item
  v = 120
\end{enumerate}

\end{exercise}

\begin{exercise}[]\protect\hypertarget{exr-line}{}\label{exr-line}

Berechne die Fahrzeit für die Strecke zum Ferienort, wenn du 1/3 der
Strecke 30 km/h fährst, 1/3 der Strecke 80 km/h fährst und 1/3 der
Strecke 100 km/h fährst.

\end{exercise}

\subsection{Lehrpersonenkommentar}\label{lehrpersonenkommentar}

\begin{tcolorbox}[enhanced jigsaw, left=2mm, colbacktitle=quarto-callout-note-color!10!white, bottomtitle=1mm, leftrule=.75mm, coltitle=black, breakable, opacitybacktitle=0.6, opacityback=0, arc=.35mm, toptitle=1mm, rightrule=.15mm, colback=white, title=\textcolor{quarto-callout-note-color}{\faInfo}\hspace{0.5em}{Hinweis}, colframe=quarto-callout-note-color-frame, titlerule=0mm, bottomrule=.15mm, toprule=.15mm]

Die Funktionen spielen eine sehr wichtige Rolle in der Mathematik.
Deshalb möchte ich für die Einführung in dieses neue Thema zwei
Lektionen beanspruchen. Mein Skript, das die Basis für diese Lektionen
bildet, beschäftigt sich vorwiegend mit Funktionsgleichungen und
Begriffserklärungen. Die SuS lernen einen neuen mathematischen Begriff,
die Funktion, sowie eine neue Notation kennen. Damit das Thema
verständlicher wird, rate ich den SuS im ersten Teil des Skripts, sich
Funktionen als kleine Maschinen vorzustellen. Dabei lernen sie bildlich,
was eine Funktion ist und wie diese funktioniert.

Im zweiten Teil des Skripts werden Funktionen anhand von Mengen
veranschaulicht. Die SuS lernen dadurch zwei verschiedene Varianten
kennen, wie man Funktionen darstellen kann. Diese Vergleiche sollen
ihnen helfen, die folgende mathematische Gleichungen besser zu
verstehen. Der Fokus dieses Skripts liegt ausführlich auf den
Funktionsgleichungen. Alternativ könnte man auch mit dem Graphen einer
Funktion beginnen und anhand dessen auf die Gleichung schliessen. Ich
halte es jedoch für sinnvoller, zuerst zu verstehen, was eine Funktion
ist und wie man sie korrekt notiert, bevor man sich mit den Arten von
Funktionen und somit ihren Graphen beschäftigt. Ich fokussiere mich
deshalb auf die Gleichungen und würde dann in einer nächsten Lektion mit
dem Graph einer Funktion fortfahren.

Nach dem ersten Teil des Skripts habe ich einige Übungen eingebaut.
Diese sind darauf ausgelegt, die SuS so intensiv wie möglich mit den
Gleichungen zu konfrontieren. Es gibt Aufgaben, bei denen sie die
Gleichung lösen müssen, und andere, bei denen sie die Funktionsgleichung
selbst formulieren sollen. Dadurch erhalten sie einen differenzierten
Überblick über das neue Thema.

Auch nach dem zweiten Teil gibt es verschiedene Aufgaben. Die erste ist
auf die Anwendung der Mengentheorie ausgelegt und beim Rest liegt der
Fokus dann wieder beim Lösen von Funktionsgleichungen. Auch habe ich ein
alltagsnahes Beispiel eingefügt, um das Thema den SuS etwas näher zu
bringen.

In diesen Lektionen liegt der Fokus auf dem neuen Thema der Funktion.
Deshalb sollten der erste und zweite Teil des Skripts zu Beginn der
Lektionen frontal eingeführt werden. So haben alle SuS die neuen
Informationen einmal gehört und es können bereits einige Beispiele
gemeinsam bearbeitet werden. Das Skript wird den SuS aber auch nach der
Einführung zugänglich sein, damit sie die Informationen zur
Ergebnissicherung noch einmal nachlesen können. Erst danach sollten sie
mit den Übungen beginnen. Als kleine ``Auszeit'' oder zusätzliche Hilfe
habe ich auch Videos eingebaut, um eine alternative Lernstrategie
anzubieten. Mir ist bewusst, dass dieser Lektionsablauf eher altmodisch
ist und man heute einen interaktiveren Unterricht planen sollte. Dennoch
bin ich überzeugt, dass es sich gerade bei der Einführung von neune
Themen sehr lohnt, wenn man sie bewusst im Frontalunterricht behandelt.

\subsubsection{Unterrichtsmethoden}\label{unterrichtsmethoden}

Um das Thema der Funktionen verständlicher zu machen, setze ich auf eine
Kombination von Unterrichtsmethoden. Zu Beginn der beiden Lektionen
führe ich das Thema frontal ein. Dabei wird den SuS das neue Thema der
Funktionen erklärt. Im ersten Teil gehe ich vor allem auf die Notation
und den Aufbau einer Funktion ein. In der zweiten Lektion wird das Thema
gewissermassen wiederholt, allerdings werden die Funktionen nun auf eine
andere Art und Weise dargestellt. Durch den Frontalunterricht erkenne
ich in der zweiten Lektion direkt, was die SuS noch nicht verstanden
haben. Die klare Einführung stellt sicher, dass alle SuS mit den neuen
Grundbegriffen arbeiten können. Gerade bei neuen Themen ist es sinnvoll,
den Frontalunterricht anzuwenden, da so alle SuS gleichzeitig auf
denselben Wissensstand gebracht werden. Anschließend kann in einer
weiteren Unterrichtssequenz auf lernschwächere oder sehr leistungsstarke
SuS individuell eingegangen werden.

Nach dem Einstieg folgt eine Phase der Einzelarbeit. Die SuS sollen das
gerade Gelernte noch einmal durchlesen, um die Informationen besser zu
verinnerlichen. Sobald dies geschehen ist, können sie beim Lösen der
Aufgaben zwischen Einzel- und Partnerarbeit wählen.

Durch die Partnerarbeit haben die SuS die Möglichkeit, sich gegenseitig
zu unterstützen. Das kooperative Lernen fördert sowohl das Verständnis
des Themas als auch ihre sozialen Kompetenzen. Außerdem haben sie durch
diese Methode die Gelegenheit, Fragen an die Lehrperson oder ihre Peers
zu stellen.

Das Austeilen des Skripts ermöglicht es den SuS, individuell an den
Aufgaben zu arbeiten. Sie können dies in ihrem eigenen Tempo tun und
dabei jederzeit die Regeln der Funktionen nachlesen, falls etwas unklar
ist. In diesen Phasen kann ich die SuS individuell unterstützen.

\subsubsection{Unterrichtstechniken Einführung
Funktionen}\label{unterrichtstechniken-einfuxfchrung-funktionen}

Bei jedem neuen Thema in der Mathematik fragt man sich, wie es für die
Schülerinnen und Schüler am besten verständlich wird. Bei den Funktionen
habe ich mich deshalb für eine visuelle Darstellung und den Einsatz von
Metaphern entschieden. Das Konzept der Funktionsgleichung vergleiche ich
mit einer kleinen Maschine, die einen Input und einen Output
verarbeitet. Diese bildhafte Beschreibung hilft den Lernenden, eine
Verbindung zwischen dem abstrakten Thema der Funktion und der Realität
herzustellen. Zudem führe ich die mathematische Notation schrittweise
ein. Zunächst arbeite ich mit der visuellen Metapher der Maschine, bevor
ich zur mathematischen Notation übergehe. Ebenso wie die Maschine helfen
auch Mengen, das Thema zu visualisieren. Da die Schülerinnen und Schüler
bereits mit Mengen vertraut sind, lassen sich einfache Alltagsbeispiele
nutzen, um das Thema verständlicher zu machen. Hier arbeite ich mit
Symbolen oder in den Übungen dann auch mit Früchten. Ich verpacke also
die symbolische Sprache der Mathematik in greifbare Beispiele. Ein
weiterer wesentlicher Bestandteil meiner Lernumgebung sind die
anschließenden Übungen, die ich nach jedem Theorieteil einfüge. Diese
Aufgaben basieren direkt auf den Beispielen aus dem Skript, sodass die
Lernenden das neue Wissen sofort anwenden können. Ich habe die Aufgaben
in Grundanforderungen und Erweiterungen unterteilt, damit alle ein
grundlegendes Verständnis erwerben, aber auch die Möglichkeit haben,
selbstständig weiterzuarbeiten und ihr Verständnis zu überprüfen. Im
Allgemeinen folgt meine Lernumgebung einer klaren und verständlichen
Struktur. Sie beginnt mit dem theoretischen Teil, gefolgt von darauf
aufbauenden Aufgaben. Was den Schülerinnen und Schülern hilft, ihr
Wissen systematisch zu vertiefen.

\subsubsection{Vertiefung 1: Lernphasen}\label{vertiefung-1-lernphasen}

Mit dem LUKAS Lernprozessmodell wird diese Lernumgebung in 5 Phasen
unterteilt.

Zu Beginn der Lektion wird der Klasse der Lernauftrag klar vermittelt.
Durch den Frontalunterricht erkläre ich das neue Thema der Funktionen.
Die SuS haben somit eine erste Vorstellung der Funktion, deren
mathematischen Notation und sie lernen eine neue Gleichungsart kennen.
Ich biete der Klasse also eine klare Orientierung, was sie lernen sollen
und warum das Thema wichtig ist. Folglich wissen die Schülerinnen und
Schüler von Anfang an, welche Lernziele sie erreichen sollen.

In der darauffolgenden Phase befinden wir uns im Bereich der
Erarbeitung. Die Klasse setzt sich also nach dem theoretischen Input mit
dem eben gelernten Stoff auseinander. Sie werden zunächst aufgefordert,
das Gelernte noch einmal zu vertiefen, indem sie es in Einzelarbeit
durchlesen. Danach können sie selbst entscheiden, ob sie die Übungen
alleine oder zu zweit erarbeiten wollen.

Die Dritte Phase, das Können, wird durch das Lösen der Übungen gestärkt.
Die SuS wenden die neuen Erkenntnisse an, indem sie die
Funktionsgleichungen lösen und auch selbständig aufstellen. Die Aufgaben
helfen den SuS ihr Verständnis über das Thema zu stärken und
gleichzeitig werden sie routinierter im Umgang mit Funktionsgleichungen.
Hierbei haben sie die Möglichkeit in ihrem eigenen Tempo zu arbeiten,
was eine individuelle Differenzierung fördert.

Die Anwendungsphase ist in dieser Lernumgebung in den letzten Aufgaben
zu finden. Mit Hilfe des Alltagsbeispiels lernen die SuS, die neuen
Informationen auf lebensnahe Kontexte zu übertragen. Dabei vertiefen sie
ihr Verständnis und kennen zugleich ein Beispiel für die Relevanz der
mathematischen Funktion im Alltag.

In der letzten Phase findet die Sicherung des Gelernten statt. Dafür
fasse ich die wichtigsten Begriffe des Themas am Ende des Skripts noch
einmal zusammen. Diese Aufzählung dient als Checkliste für die
Schülerinnen und Schüler, falls sie eine Begriff nicht mehr definieren
können. Auch dient sie als Anhaltspunkt für die Lehrperson beim
Überprüfen, welche Konzepte noch nicht erfasst werden konnten.

\subsubsection{Vertiefung 2: Motivation}\label{vertiefung-2-motivation}

Die Motivation der Schülerinnen und Schüler spielt in jeder Lernumgebung
eine zentrale Rolle. In den Mathematikunterricht der Oberstufe kommen
viele SuS leider bereits mit negativen Erfahrungen oder sogar Ängsten
daher. Deshalb ist es umso wichtiger, ihre Motivation anzuzapfen und
stetig aufzubauen.

In dieser Lernumgebung setze ich deshalb auf eine klare Struktur und
verständliche Aufgabenstellungen, um die Ängste der SuS zu minimieren.
Die frontale Einführung gibt ihnen eine solide Grundlage, auf der sie im
weiteren Verlauf der Lektion aufbauen können. Falls der Einstieg für
manche zu schnell war oder sich einzelne SuS nicht genug konzentrieren
konnten, stellt das Skript einen weiteren Meilenstein dar. Es bietet
ihnen die Möglichkeit, immer wieder auf die Theorie zurückzugreifen,
falls sie auf Schwierigkeiten stossen. Zusätzlich steht ihnen ein
unterstützendes Video zur Verfügung. Dieses Video ist nicht offizieller
Teil der Lernumgebung, dient jedoch als zusätzliche differenzierte
Hilfe. Sobald die SuS das neue Thema verstanden haben, steigt auch ihre
Motivation. Denn am wenigsten motiviert ist man, wenn man einen
Themenbereich nicht nachvollziehen kann. Durch verschiedene
Unterrichtsmethoden und gezielt eingesetzte Techniken wird die
Motivation der SuS weiter gestärkt. Sie haben die Freiheit, ihr
Lerntempo selbst zu gestalten und können zwischen Einzelarbeit und
kooperativem Lernen wählen. Bildhafte Darstellungen und Metaphern helfen
ebenfalls, mögliche Hemmungen abzubauen und das Thema zugänglicher zu
machen. Ein weiterer wichtiger Punkt ist das Einbeziehen von
Alltagsbeispielen. Diese helfen den SuS, die Relevanz komplexer
mathematischer Themen in ihrem täglichen Leben zu erkennen, sodass das
Thema nicht als „unnötig`` wahrgenommen wird. Zusammenfassend tragen
alle diese Faktoren dazu bei, dass die Schülerinnen und Schüler
Interesse am Thema entwickeln, ihre Aufgaben motiviert angehen und das
Gelernte langfristig festigen.

\end{tcolorbox}



\end{document}
